\chapter{Marco práctico: el caso de estudio SMCI}
\label{cap:marco-practico}

% Cifras desde RESULTADOS_OBJETIVO.md §1bis (SMCI) / outputs/experiments/m10_smci_*.json.
% La tabla principal debe autogenerarse; aquí van los valores canónicos del §1bis.
% Figuras: se generan desde notebooks/STRATA_SMCI.ipynb (ver graficas_clave.md). Placeholders marcados.

Este capítulo responde a la pregunta operativa del trabajo sobre un activo concreto: en un benchmark justo, ¿un meta-learner desplegable que consume las señales de STRATA mejora la dirección del agente y supera a las estrategias triviales? La respuesta sobre SMCI es afirmativa en términos nominales y robusta a varias particiones de los datos, con la significancia plena delimitada como trabajo futuro. El capítulo desarrolla esa respuesta en orden: los datos y por qué SMCI es un banco de pruebas honesto; el agente sin supervisar; la regla a mano; el meta-learner M10 y su resultado; la robustez; la interpretabilidad por features; y una lectura honesta de los límites.

\section{Datos y caso de estudio}
\label{sec:datos}

El activo del caso de estudio es \textbf{SMCI}. La elección responde a un criterio metodológico: sus dos clases direccionales están casi equilibradas, con un \emph{buy-and-hold} (comprar y mantener, siempre largo) que acierta el $48{,}4\,\%$ de los días. En un activo con clases balanceadas la accuracy es una métrica informativa y la comparación con las estrategias triviales es exigente, porque la regla de no-habilidad (predecir siempre la clase mayoritaria, \emph{ZeroR} en la terminología de aprendizaje automático \cite{witten2016datamining, kuhn2008caret}) no se lleva por delante a cualquier modelo solo por el desbalance. En SMCI esa regla mayoritaria acierta el $51{,}6\,\%$, materializada como ``siempre corto'', porque bajan algo más días de los que suben.

Los tres detectores de STRATA se calibran una sola vez sobre 2000--2024-09 y se congelan, según la Sección~\ref{sec:strata-calib}. La evaluación se hace fuera de muestra sobre los días posteriores, de los que se descartan los primeros como periodo de arranque del walk-forward, lo que deja $n = 250$ días puntuables. El inicio del periodo de evaluación es posterior al corte de conocimiento del modelo de lenguaje que mueve al agente, de modo que las decisiones del agente sobre estos días no pudieron memorizarse durante su entrenamiento. La causalidad se garantiza con la convención \texttt{signal\_lag=1} y con el protocolo walk-forward de embargo $=1$ de las Secciones~\ref{sec:cpcv-causalidad} y~\ref{sec:cpcv-walkforward}.

\section{El agente sin supervisar (M5)}
\label{sec:m5}

El agente de trading por sí solo, que denotamos M5, es un perdedor direccional sobre SMCI: acierta el $48{,}4\,\%$ de las direcciones, por debajo del azar. Su comportamiento está además muy sesgado, con posición corta el $95\,\%$ de los días. Este es el punto de partida del problema: un decisor que, sin supervisar, pierde dinero y no alcanza el medio acierto. La hipótesis del trabajo es que las señales de STRATA, bien explotadas, recuperan parte de esa dirección perdida.

\section{La regla a mano (M8) como referencia}
\label{sec:m8}

La regla determinista M8 aplica la intervención \emph{override-C} de la Sección~\ref{sec:strata-intervencion}: cuando el RAM score supera el gate $\tau = 0{,}50$, sustituye la dirección del agente por la que marca el régimen. Sobre SMCI, M8 acierta el $49{,}6\,\%$, apenas por encima de M5. La razón es mecánica y se reporta con transparencia: el agente ya va casi siempre corto, y el régimen de SMCI también sesga a corto, de modo que la dirección del agente y la del régimen coinciden la mayor parte del tiempo y STRATA solo encuentra incoherencia que corregir en torno al $3\,\%$ de los días. M8 cumple aquí el papel de \emph{referencia interpretable}: una regla transparente sobre la que medir cuánto añade un modelo que aprende.

\section{El meta-learner M10}
\label{sec:m10}

M10 es el modelo central del caso de estudio. Es un clasificador binario XGBoost \cite{chen2016} que, para cada día, estima la probabilidad de que el retorno del día siguiente sea positivo a partir de \textbf{veintidós variables}: las quince de la decisión del agente (cinco personalidades, cada una con signo, tamaño y confianza) y las siete señales de STRATA y de régimen (los tres scores RAM, PSA y GSO, las tres probabilidades de régimen y la volatilidad condicional del GARCH). La posición del día es el signo de la probabilidad estimada menos $0{,}5$.

El modelo se valida y se opera con el protocolo walk-forward de la Sección~\ref{sec:cpcv-walkforward}: un periodo de arranque de $150$ días, reentrenamiento cada $21$ días con toda la información pasada, y embargo de un día. Para reducir la varianza del azar interno del algoritmo se promedia un \emph{ensemble} de diez modelos que solo difieren en la semilla, técnica de \emph{bagging} respaldada por la literatura \cite{breiman1996, dietterich2000}; el promedio cancela ruido sin crear señal, así que no introduce optimismo. M10 no usa la validación cruzada combinatoria (CPCV), que mira el futuro y se reserva como contraste (Sección~\ref{sec:cpcv-cpcv}); tampoco incorpora variables de momento ni un mecanismo de abstención, descartados porque no mejoran la accuracy fuera de muestra.

\section{Resultado principal}
\label{sec:resultado}

La Tabla~\ref{tab:smci-resultado} resume el resultado sobre todo el periodo de evaluación. \textbf{M10 acierta el $55{,}2\,\%$ de las direcciones, por encima del agente ($48{,}4\,\%$), de la regla a mano ($49{,}6\,\%$), de comprar-y-mantener ($48{,}4\,\%$) y de la clase mayoritaria ($51{,}6\,\%$).} Es el único de los seis competidores que supera a la vez a los tres baselines y a las dos estrategias triviales. En el plano económico, ilustrativo y no probatorio, M10 multiplica el capital por $3{,}24$ frente al $0{,}71$ de comprar-y-mantener, con un ratio de Sharpe de $+1{,}84$.

\begin{table}[htbp]
\centering
\caption{Resultado sobre SMCI (OOS, $n=250$ días, walk-forward, embargo$=1$). Fuente: \texttt{m10\_smci\_valtest\_robustez.json} / \texttt{RESULTADOS\_OBJETIVO.md} §1bis.}
\label{tab:smci-resultado}
\begin{tabular}{lrrr}
\hline
Estrategia & Accuracy & Sharpe & Equity \\
\hline
\textbf{M10-WF ensemble (modelo)} & \textbf{0,552} & \textbf{+1,84} & \textbf{3,24$\times$} \\
Trivial: clase mayoritaria (= siempre corto) & 0,516 & --- & --- \\
M8 (regla a mano STRATA) & 0,496 & +0,33 & 1,02$\times$ \\
M5 (agente solo) & 0,484 & $-$0,24 & 0,98$\times$ \\
Trivial: comprar y mantener (siempre largo) & 0,484 & +0,03 & 0,71$\times$ \\
\hline
\end{tabular}
\end{table}

% [Figura por generar — graficas_clave §1: barras de accuracy M10 vs M5/M8/B&H con líneas de azar (0,5) y mayoritaria (0,516); + curvas de equity. Fuente: notebooks/STRATA_SMCI.ipynb.]

La honestidad estadística exige acompañar el $0{,}552$ de su contraste. La ventaja de M10 es \textbf{nominal, no significativa} al tamaño de muestra disponible: el contraste binomial frente a la clase mayoritaria, que es el baseline de no-habilidad correcto, da $p = 0{,}14$; el test de permutación por bloques frente a comprar-y-mantener da $p = 0{,}047$, que no sobrevive a la corrección por multiplicidad; y el Deflated Sharpe Ratio queda en $0{,}72$, por debajo del umbral de habilidad. Esto es lo esperable y no contradice la tesis. Con un horizonte de evaluación de unos $250$ días, acotado porque el agente solo existe sobre el periodo posterior al corte del modelo de lenguaje, ningún sistema de decisión diaria alcanza significancia plena frente a un benchmark justo; es una consecuencia conocida del tamaño de muestra y de la eficiencia de los mercados a horizonte diario. El criterio del trabajo, por tanto, es el que la literatura y la práctica admiten en este régimen de muestra: superar a lo trivial de forma nominal y \emph{consistente}, y dejar la significancia plena como trabajo futuro con un periodo más largo. La consistencia es lo que se establece a continuación.

\section{Robustez}
\label{sec:robustez}

La ventaja de M10 no depende del corte concreto de los datos. La Tabla~\ref{tab:smci-robustez} la evalúa sobre tres particiones cronológicas estándar, $60/40$, $70/30$ y $80/20$: en las tres, y tanto en el tramo de validación como en el de test, M10 supera al agente, a la regla y a comprar-y-mantener. La conclusión es invariante a dónde se corte la serie, lo que descarta que el resultado sea un artefacto de una partición afortunada.

% [Tabla/figura por generar — graficas_clave §3: robustez a la partición (3 splits val/test). Fuente: m10_smci_valtest_robustez.json / notebooks/STRATA_SMCI.ipynb.]

\begin{table}[htbp]
\centering
\caption{Robustez a la partición: accuracy de M10 frente a los baselines en validación y test de tres cortes cronológicos. M10 gana a todos en los tres. Fuente: \texttt{m10\_smci\_valtest\_robustez.json}.}
\label{tab:smci-robustez}
\begin{tabular}{lcc}
\hline
Partición (val / test) & Accuracy M10 (val) & Accuracy M10 (test) \\
\hline
60 / 40 & 0,52 & 0,60--0,62 \\
70 / 30 & 0,52--0,535 & 0,60--0,62 \\
80 / 20 & 0,52--0,535 & 0,60--0,62 \\
\hline
\end{tabular}
\end{table}

La robustez se confirma en el tiempo. Sobre ventanas rodantes a lo largo del periodo de evaluación, M10 supera a comprar-y-mantener en el $71$--$82\,\%$ de las ventanas y al agente en el $67$--$76\,\%$, según el tamaño de ventana. Que la ventaja se sostenga en la mayoría de sub-periodos, y no en uno solo, es la evidencia de que no se trata de la suerte de un tramo concreto. La elección del protocolo también es estable: la accuracy apenas se mueve al variar la frecuencia de reentreno, y el embargo se fija en uno por el horizonte de la etiqueta (Sección~\ref{sec:cpcv-purga}), no para maximizar una cifra.

% [Figura por generar — graficas_clave §3/§5: rolling-window (% de ventanas ganadas) y robustez al embargo. Fuente: notebooks/STRATA_SMCI.ipynb §G, §0bis.]

\section{Interpretabilidad: M10 es inseparable de STRATA}
\label{sec:interpretabilidad}

Queda justificar que la ventaja de M10 proviene de las señales de STRATA y no de cualquier conjunto de variables. Dos análisis lo establecen. El primero es una ablación: un M10 entrenado solo con las quince variables del agente, sin las siete señales de STRATA, cae al $46{,}8\,\%$, el nivel del agente; añadir las siete señales lo lleva al $55{,}2\,\%$, una mejora pareada con $p = 0{,}05$ en el test de McNemar. El segundo es el análisis de importancia por valores de Shapley (SHAP) \cite{lundberg2017} sobre el ensemble: las siete señales de STRATA y de régimen concentran el $41{,}4\,\%$ de la importancia total, por delante de las variables del agente. Sin las señales de STRATA, M10 no añade nada sobre el agente; con ellas, recupera la dirección. El meta-learner es, por construcción, inseparable de STRATA, lo que da contenido al nombre del sistema.

% [Figura por generar — graficas_clave apéndice: ablación (agente-solo vs +STRATA) y barras SHAP (7 señales STRATA arriba). Fuente: notebooks/STRATA_SMCI.ipynb §7, §7b.]

\section{Lectura honesta y límites}
\label{sec:lectura}

El resultado tiene un alcance acotado. En SMCI los tres modelos, M5, M8 y M10, apuestan en la misma dirección la mayor parte del tiempo, porque el agente ya está alineado con un régimen que también sesga a corto. La supervisión de STRATA rescata de forma marcada al agente solo cuando este va a contracorriente de un régimen que acierta, una condición que SMCI cumple poco y que delimita el alcance del método. Aun así, sobre el periodo completo M10 mantiene una posición corta en torno al $47\,\%$ de los días, de modo que su acierto no se reduce a apostar siempre a la baja: gana a la estrategia constante ``siempre corto''.

La limitación principal es el tamaño y la naturaleza de la muestra. El agente solo existe sobre el periodo posterior al corte del modelo de lenguaje, lo que acota la evaluación a unos $250$ días y vuelve inalcanzable la significancia plena frente a un benchmark justo. Las dos extensiones inmediatas, la generalización a otros agentes y a otros activos, son costosas en inferencia y en tiempo por las razones que se cuantifican en el Capítulo~\ref{cap:conclusiones}, y quedan como trabajo futuro. La aportación que sí se sostiene es la de un protocolo de supervisión estadística interpretable y desplegable que, en un benchmark honesto, recupera la dirección del agente y supera a la regla a mano y a lo trivial.
